\section{Introduction}
DC motors are fundamental components in control engineering, valued for their simplicity, 
reliability, and adaptability across a wide range of applications. They are commonly used in 
systems where precise control of speed and torque is essential, making them ideal for diverse 
industries and applications.

In everyday applications, DC motors are integral to automotive systems, including electric 
power steering and drive mechanisms for various car components. In the medical field, they 
drive devices that require smooth and reliable operation, such as patient beds and ventilators. 
Industrial automation heavily relies on DC motors to enable precise movement in robotic arms, 
conveyor systems, and automated assembly lines. In the entertainment industry, DC motors power 
animatronics and stage effects, enhancing performance and visual impact. Home appliances such as 
fans, vacuum cleaners, and washing machines also use DC motors for efficient energy conversion. 
Finally, in aerospace, DC motors contribute to systems that require dependable operation under 
extreme conditions, from instrumentation to control surfaces in aircraft.

This lab explores the fundamental principles governing DC motor behavior and their practical 
applications, highlighting the role of DC drives in enabling controlled motion and efficient 
energy use across different fields.

\textbf{Learning Goals}\\
This lab aimed to deepen our understanding of DC motor control, specifically by exploring methods to 
measure, model, and regulate motor behavior in a controlled manner. The primary goals were to develop 
skills in approximating transfer functions from measured responses, to model and simulate motor dynamics 
within MATLAB’s Simulink environment, and to design an effective PID controller for the system. 
By iteratively testing and refining PID values, we sought to achieve a stable, responsive control 
setup that could be applied to a real motor system. Additionally, this lab emphasized the practical 
application of control theories, such as root locus analysis, and allowed us to experience firsthand 
how tuning and testing control parameters impact system performance.\\
\\
\textbf{Motivation}\\
This exercise was motivated by the practical need for control engineers to design efficient and responsive 
systems across various industries, where accurate motor control is essential. Gaining familiarity with modeling 
and controller design enhances our capabilities in fields such as robotics, automotive systems, and industrial 
automation, where motor control plays a critical role.
\newpage
\section{Procedure and Method}

The lab was structured to progress through a series of steps that systematically built our knowledge and 
skills in DC motor control. 

Before the hands-on portion of the laboratory, several essential preparatory tasks were completed. 
esearch was conducted on real-life applications of DC drives across various industries, including 
automotive, medical, industrial automation, entertainment, home appliances, and aerospace. 
This research established context by illustrating the broader relevance of DC motor control 
and how similar methodologies could be applied across diverse fields.

The derivation of the transfer functions for a DC motor followed, focusing on responses related to both 
angular velocity and rotational speed.

This process utilized the \grqq white-box modeling\grqq approach. However, applying this approach requires 
knowledge of the motor's parameters, which were unavailable for the laboratory setup. Instead, MATLAB was 
used to perform a system-identification task by fitting a PT2 system function to a measured step response, 
with results detailed in subsequent sections.

With the motor’s transfer function established, the system was implemented in Simulink to simulate its 
performance. This simulation enabled analysis of the system’s response and facilitated the design and 
inclusion of a PID controller. Root locus analysis in MATLAB was applied to examine the movement of poles 
under varying gain conditions, identifying the critical value for stability. Based on this, the Ziegler-Nichols 
method was employed to calculate initial values for the PID controller, which were then iteratively refined to 
optimize the motor’s response.

The refined controller was subsequently tested on the physical DC motor. This practical application 
highlighted discrepancies between simulated and real-world responses, emphasizing the importance of 
tuning and parameter adjustments to achieve alignment. This approach led to the development of a reliable 
control setup that provided stable, responsive performance.

The methodology reinforced foundational knowledge in DC motor modeling and control design while also 
developing practical skills in root locus analysis, gain tuning, and the application of the Ziegler-Nichols method.



\newpage
\section{Results}

\subsection{Preparatory work}

The solution to the preparatory work begins with the derivation of the transfer function of a DC motor, 
utilizing Kirchhoff's first law and the equations of motion. The resulting expression for the angular 
velocity is shown in Equation (\ref{eq_angular_v_tf}), where $U_{in}$ represents the applied voltage at the 
terminals, $\omega$ denotes the angular velocity, $J_{mot}$ is the inertia of the rotor, $b_{mot}$ is the 
damping coefficient of the rotor, $L_{mot}$ is the inductance of the coil, $R$ is the resistance of the 
wires, $k_e$ is the back electromotive force constant, and $k_t$ is the torque constant.

\begin{equation}
    \mathcal{G}_{\omega}(s)=\frac{\omega (s)}{U_{in}(s)} = \frac{k_t}{(J_{mot}*s+b_{mot})*(L_{mot}*s+R_{mot})+k_e*k_t}
    \label{eq_angular_v_tf}
\end{equation}

By multiplying equation (\ref{eq_angular_v_tf}) with an integrator, it is possible to find the corresponding 
transfer function for the angular displacement $\Theta$, which is done in equation (\ref{eq_angular_disp_tf}).

\begin{equation}
    \mathcal{G}_{\Theta}(s) = \frac{\Theta (s)}{U_{in}(s)} = \frac{k_t}{s*((J_{mot}*s+b_{mot})*(L_{mot}*s+R_{mot})+k_e*k_t)}
    \label{eq_angular_disp_tf}
\end{equation}

The researched applications of DC motors in real-life scenarios highlight their versatility and significance 
across various industries. In the automotive sector, DC motors are utilized in electric power steering systems, 
propulsion mechanisms for electric vehicles, and other precision-driven components. In the medical field, these 
motors drive devices such as ventilators, patient beds, and advanced surgical equipment. Notably, the Da Vinci 
surgical robot, known for its precision in minimally invasive surgeries, employs the same type of DC motors used 
in the laboratory task, underscoring the real-world relevance of the experimental setup.

DC motors are also extensively applied in industrial automation, powering robotic arms, conveyor systems, and 
CNC machines. In the entertainment industry, they enable smooth motion for animatronics and stage effects, while 
in household appliances, they contribute to efficient energy conversion in devices like vacuum cleaners and fans. 
The aerospace sector leverages DC motors for lightweight and reliable actuators in aircraft systems and 
instrumentation under extreme conditions.

Additionally, measurement data was provided, including the step response of a DC motor, to facilitate system 
identification. The laboratory setup consisted of a preconfigured system comprising a DC motor and an encoder 
device for measuring angular position. This configuration enabled a detailed analysis of the motor’s behavior 
and performance in controlled conditions, bridging the theoretical insights with practical implementation.

\subsection{Identification of the system}

The provided Simulink model was executed to measure the output velocity of the DC motor in response to a step 
input of 12\,\si{\volt} over a duration of 1.5 seconds. This process generated a table of measured data, which 
was subsequently plotted and visualized using MATLAB to facilitate further analysis. From the collected data, a 
transfer function was derived by fitting a PT2 system to the observed step response. The resulting transfer 
function is presented in Equation (\ref{eq_fitted_tf_omega}).

\begin{equation}
    \mathcal{G}_{\omega} = \dfrac{275.9}{s^2+34.06*s+394.6}
    \label{eq_fitted_tf_omega}
\end{equation}

The fitted step response of the system regarding the angular velocity $\omega$ can be viewed in 
figure \ref{fig:fitted curve}, along with the measurements.
\begin{figure}[H]
    \centering
    \includegraphics[width=0.9\linewidth]{Images/sys_resp_omega.pdf}
    \caption{step response of fitted curve over measured data (angular velocity)}
    \label{fig:fitted curve}
\end{figure}

By multiplying $\mathcal{G_{\omega}}$ with an integrator, the transfer function for the angular displacement 
comes out as $\mathcal{G}_{\Theta}$, which can be seen in equation (\ref{eq_fitted_tf_theta}).

\begin{equation}
    \mathcal{G}_{\Theta} = \dfrac{275.9}{s^2+34.06*s+394.6} * \dfrac{1}{s} = \dfrac{275.9}{s^3+34.06*s^2+394.6*s}
    \label{eq_fitted_tf_theta}
\end{equation}


\newpage
\subsection{Controller design}

Using root-locus analysis in MATLAB, it was possible to determine the critical gain $K_{crit}$, which could 
then be used as a proportional controller in a feedback loop with $\mathcal{G}_{\Theta}$ to find the period 
time of the oscillating system $T_{osc}$. This is visualised in figure \ref{fig:sys_theta_oscillating}.


\begin{figure}[H]
    \centering
    \includegraphics[width=0.9\linewidth]{Images/sys_resp_theta_osc.pdf}
    \caption{Step response of $\mathcal{G}_{\Theta}$ with proportional controller using critical gain $K_{crit}$, 
    visualising $T_{osc}$}
    \label{fig:sys_theta_oscillating}
\end{figure}


Knowing $K_{crit}$ and $T_{osc}$, it was possible to use the Ziegler-Nichols method to calculate the 
values $K_p$, $K_i$ and $K_d$ as a starting point for a PID-controller using table \ref{tab:ziegler_nichols_2}.

\newline
\begin{table}[H]
\centering
\caption{Ziegler-Nichols table for determining starting values for P-, PI- or PID-controllers}
\label{tab:ziegler_nichols_2}
\begin{tabular}{|l|lll|}
\hline
Controller & \multicolumn{3}{l|}{Parameters}                                                                                    \\ \hline
P          & \multicolumn{1}{l|}{$K_p=0.5*K_{crit}$}  & \multicolumn{1}{l|}{}                             &                     \\ \hline
PI         & \multicolumn{1}{l|}{$K_p=0.45*K_{crit}$} & \multicolumn{1}{l|}{$T_i=0.83*T_{osc}$} &                     \\ \hline
PID        & \multicolumn{1}{l|}{$K_p=0.6*K_{crit}$}  & \multicolumn{1}{l|}{$T_i=0.5*T_{osc}$}            & $T_d=0.125*T_{osc}$ \\ \hline
\end{tabular}
\end{table}


Substituting the values of our system in equations (\ref{eq_ziegler_nichols}), being $K_{crit} = 48.7$ and $T_{osc} = 0.316$, we get:
\begin{subequations}
    \begin{equation}
        K_p = 0.6*K_{crit} = 0.6 * 48.7 = 29.22
    \end{equation}
    \begin{equation}
        K_i = \dfrac{K_p}{T_i} = \dfrac{K_p}{0.5*T_{osc}} = \dfrac{29.22}{0.5 * 0.316} = 184.77
    \end{equation}
    \begin{equation}
        K_d = K_p * T_d = K_p*0.125*T_{osc} = 29.22*0.125 * 0.316 = 1.155
    \end{equation}
    \label{eq_ziegler_nichols}
\end{subequations}

Inserting these values into a PID controller and testing our controller on our identified transfer 
function $\mathcal{G}_{\Theta}$ in a feedback loop using Simulink, we measured the step response of our 
new system, which can be seen in figure \ref{fig:ZNCurve}.

\begin{figure}[H]
    \centering
    \includegraphics[width=0.9\linewidth]{Images/sys_resp_first_pid.pdf}
    \caption{Simulated step response of the system with PID controller using Ziegler-Nichols parameters to a 
    step input of 30\si{\degree}}
    \label{fig:ZNCurve}
\end{figure}
While far from perfect, these values for the controller were a good starting point for further tuning. 
We were now ready to apply them to the real system.

\newpage

\subsection{Controller implementation and evaluation}
During the laboratory tasks, issues arose with the PID parameters, preventing the direct application of values 
determined using the Ziegler-Nichols method to the real system. Instead, the PID-Tuner applet in MATLAB was 
utilized to generate starting values for the controller parameters. The identified error was traced to the 
calculation of the $K_d$ value. However, due to these challenges, no measurements on the real system are available 
for the parameters derived from the Ziegler-Nichols method.

\begin{figure}[H]
    \centering
    \includegraphics[width=1.0\linewidth]{Images/30deg_real.pdf}
    \caption{Response of real system when exposed to a step input of \ang{30} for different PID-parameters}
    \label{fig:real_sys_30_deg}
\end{figure}

As shown in Figure \ref{fig:real_sys_30_deg}, altering the parameters of the controller has a significant 
impact on the system's response. To determine suitable values for the PID parameters, an iterative approach 
was employed.

Initially, the PID Tuner Toolbox was utilized to obtain approximate values for the proportional ($K_p$), 
integral ($K_i$), and derivative ($K_d$) gains. These values served as a starting point for further refinement. 
The process began with adjustments to $K_p$ and $K_i$, while $K_d$ was temporarily fixed. By iteratively 
varying $K_p$ and $K_i$, an improved response was achieved.

Subsequently, $K_d$ was introduced into the tuning process, allowing further refinement of the system's dynamic 
behavior. This iterative procedure ensured that the controller parameters were optimized to achieve the desired 
performance, balancing stability and responsiveness. For interest's sake, we also tried what would happen to the 
system, if inappropriately high values for the gains were chosen, which are also visualised in the lower sub-plots 
of figure \ref{fig:real_sys_30_deg}. The system begins to oscillate when being controlled using these values.

\begin{figure}[H]
    \centering
    \includegraphics[width=1.0\linewidth]{Images/30sine_real.pdf}
    \caption{Response of real system when exposed to a sine wave input with an amplitude of \ang{30} for different 
    PID-parameters}
    \label{fig:real_sys_30_sine}
\end{figure}

The system's response was further evaluated using a sinusoidal input instead of a simple step input, allowing for 
the observation of its behavior under various PID parameter configurations. The results of these experiments are 
presented in Figure \ref{fig:real_sys_30_sine}. Similar to the step input scenario, an excessive increase in the 
proportional gain ($K_p$) resulted in significant oscillations in the system's output. This underscores the 
importance of precise tuning of $K_p$ to achieve a balance between stability and responsiveness.

\newpage
\section{Discussion and Interpretation}

The results of this laboratory exercise offer valuable insights into the design, simulation, and practical 
implementation of DC motor control systems. Through methodical analysis and testing, we were able to explore 
the motor’s dynamics, design a control strategy, and implement it effectively.
\subsection{Transfer Function and System Behavior}

The derivation of the transfer functions for angular displacement and velocity was a foundational step in 
understanding the motor's behavior. By approximating the motor's response through curve fitting, we achieved a 
mathematical representation that closely mirrored the real system. This demonstrates the utility of such 
approximations in simplifying complex systems without sacrificing critical accuracy. The transfer function 
provided a basis for subsequent simulations and analyses, allowing us to predict system behavior and design a 
controller tailored to its characteristics.

\subsection{Root-Locus Analysis and Stability}

The root-locus analysis proved instrumental in evaluating the system’s stability and the effects of varying 
controller gain. By mapping the system poles and their movement with changes in gain, we identified the critical 
gain—an essential parameter that marks the threshold of stability. This visualization is crucial in control 
design, as it helps ensure that the system operates within a safe and effective range. The clarity provided by 
the root-locus approach also underscores its importance as a diagnostic tool in control engineering.

\subsection{The Role of Ziegler-Nichols in PID Design}

The Ziegler-Nichols method plays a pivotal role in determining the initial values for the proportional ($K_p$), 
integral ($K_i$), and derivative ($K_d$) gains of the PID controller. This empirical tuning method is particularly 
valuable because it provides a systematic approach to achieve a responsive and stable system. By starting with the 
critical gain and oscillation period, the Ziegler-Nichols rules allows to derive parameters that will result in a 
controller capable of balancing responsiveness and stability effectively. Due to aforementioned calculative errors,
this step was replaced by utilization of the PID-Tuner Toolbox within MATLAB to allow taking further steps through 
the lab. This Toolbox utilizes the Ziegler-Nichols approach, as well as it employs other PID-Tuning algorithms 
that are not tought in this lab. Finally further iterative refinement was necessary to tailor the controller to 
our specific setup, the method gave a strong starting point, reducing the time and effort required for tuning.

The refinement process, carried out in Simulink, allowed us to fine-tune the controller to achieve desired 
performance characteristics. Simulations revealed the impact of each PID component: the proportional term 
improving response speed, the integral term reducing steady-state error, and the derivative term dampening 
oscillations. These insights are critical for understanding how control parameters influence system behavior.

\subsection{Simulated and Real-World Performance}

The comparison between simulated results and the real-world motor system highlighted both the strengths and 
challenges of practical control design. The PID controller performed well in simulation, demonstrating smooth 
and stable responses. However, minor discrepancies arose during real-world testing due to factors such as 
mechanical friction, electrical noise, and sensor inaccuracies, which were not accounted for in the model. 
This underscores the importance of real-world testing as a complementary step to simulation, bridging the gap 
between theoretical predictions and practical performance.

\subsection{Practical Implications}

The insights gained through this lab have direct relevance to a wide array of applications where precise motor 
control is essential. For example, in industrial automation, the ability to design robust and responsive 
controllers ensures the smooth operation of robotic arms and conveyor systems. Similarly, in automotive systems, 
fine-tuned motor control is critical for electric power steering and other precision-driven mechanisms. 
The exploration of Ziegler-Nichols tuning and root-locus analysis equips engineers with powerful tools to design 
and optimize such systems efficiently.

\subsection{Source code}
For a short description on what each of the files in the source code does, please consider having a look at 
README.md.

